\subsection{Both sign cosets contain actual nonpentagonal elements}

\begin{statement}{PC.1}{Proved}{The kernel of the sign is strictly larger than its pentagon subgroup}
Under the supplied R1, R2 and R4, both fibres of
$\Delta:\GZ\to\Ftwo$ contain actual multiplier-one elements
which fail the profinite pentagon. In particular,
\[
 G_{\Q}\ \subsetneq\ \ker\Delta,
 \qquad
 \GT\cap\ker\Delta\ \subsetneq\ \ker\Delta.
\]
If the main vanishing assertion holds, its inclusion chain is
therefore necessarily
\[
 G_{\Q}\ \subseteq\ \GT\ \subsetneq\ \ker\Delta
                       \ \subsetneq\ \GZ.
\]
Only the first inclusion in that conditional chain remains of
unknown strictness here. This statement does not decide the
main vanishing assertion.
\end{statement}
\begin{proof}
Fix $\ell\ge5$ and let $g\in\HH$ be the actual R1 lift in
LP.5 of the independently constructed nonpentagonal
pro-$\ell$ pair. Let $k\in\HH$ be the R2 element, with
$\Delta(k)=1$ and identity action on
$\widehat B/R_{\rm nil}(U)$, where $U=F'F^3$.
For every homomorphism from $F$ to a finite $\ell$-group,
its restriction to $U$ has nilpotent image. Thus
\[
 R_{\rm nil}(U)\subseteq O^\ell(F).
\]
The action of $k$ on the actual relative pro-$\ell$ quotient
$\widehat B/O^\ell(F)$ is consequently identity.
Its normalized pro-$\ell$ parameter is one: it centralizes
$y$, belongs to $P_\ell'$, and
$C_{P_\ell}(y)=\langle y\rangle$ has trivial intersection
with that derived subgroup. Hence multiplying $g$ on either
side by any power of $k$ retains the complete normalized
pro-$\ell$ pair, not just its leading term.

Write $e=\Delta(g)\in\{0,1\}$ and set
$h_0=gk^e$, $h_1=h_0k$. Both are actual elements of $\HH$.
The character law gives $\Delta(h_i)=i$, while their
pro-$\ell$ pentagon defects retain the nonzero coefficient
$16$ from LP.4. A profinite pentagon solution maps to a
pro-$\ell$ pentagon solution by the actual braid quotient
used there, so neither $h_i$ is in $\GT$.
In particular $h_0\in\ker\Delta\setminus\GT$ and is not
Galois. This proves both strict inclusions. If
$\Delta|_{\GT}=0$, replace $\GT\cap\ker\Delta$ by $\GT$;
the final strict inclusion follows from $\Delta(k)=1$.
The proof needs no guessed value of $\Delta(g)$: the two
case-dependent definitions select the required element
using the already defined finite continuous character.
\end{proof}
